% Owner: Hoang. Interpret only results that have been measured and independently checked.
\section{Discussion and Limitations}\label{sec:discussion}

\subsection{Interpretation of Results}
The planned experiments can support different conclusions depending on what they reveal. A low API latency under light load would establish only that the tested configuration handled that load; it would not establish scalability. A successful replay test would support idempotency for the tested event and consumer version, not for every consumer. A recovery trace showing an outbox row eventually published would support the specific failure path exercised, while other crash points could still remain untested. The discussion should therefore pair each claim with its workload, configuration, sample size, and raw evidence. \todo{Interpret actual measurements and report confidence intervals or run-to-run variation where appropriate.}

The most informative result may be a trade-off rather than an improvement. For example, an event-driven path may shorten an initiating request while extending the time until all downstream views agree. The relevant service-level requirement determines whether that trade-off is acceptable. The paper should state which operation completes at the API response and which effects are expected later, then compare observed delays with a justified requirement. Until that requirement and those measurements exist, no statement of adequacy is warranted.

\subsection{Architectural Trade-offs}
Separating services can create clearer ownership and independent release boundaries, but it adds network calls, operational components, and eventual consistency. A database-per-service rule prevents accidental shared writes only if service interfaces and access controls enforce it. OpenFeign can make dependencies explicit, yet a long synchronous chain can transmit failure and latency to the user. Kafka can decouple downstream processing, yet it requires schema governance, consumer monitoring, replay controls, and a defined response to messages that repeatedly fail.

The transactional outbox reduces the risk that a local commit has no corresponding publication intent. It does not remove repeated delivery, so consumer-side idempotency remains necessary. Idempotency itself depends on stable identifiers and a transactional relationship between the deduplication record and the business effect. These mechanisms should be evaluated together: measuring only broker delivery would miss an incorrectly duplicated billing record. The design should also examine operational cost and complexity once deployment evidence is available. \todo{Compare measured costs of the chosen mechanisms with an explicit baseline or justified requirements.}

\subsection{Threats to Validity and Limitations}
The principal current limitation is missing source evidence. The system specification was not present in the workspace, so the service descriptions, workflows, security rules, and event table are provisional. A later source-code audit may reveal fewer implemented services, different data boundaries, or no outbox relay. The paper must be corrected rather than treating the current design as a historical fact.

Internal validity can be weakened by changing configuration between trials, unobserved background traffic, warm caches, or fault injections that occur at a different point than intended. Construct validity can be weakened when an event appearing in Kafka is counted as a completed business effect, or when API success hides a later failure. External validity is limited if tests use one machine, synthetic data, or a workload unlike actual rental operations. The final paper should document these conditions and avoid extrapolating to production deployments or other rental platforms. The literature review is also incomplete: the two foundational books establish patterns and semantics, but do not establish a rental-domain research gap. \todo{Update this section after the specification audit, literature review, and experiments.}
